\section{A proof of the period-function lemma}
\label{app:period-torsor}

The existence of the functions $\tau,\mu,\beta$ with the required
transformation laws is proved in \citep[Theorem~3.4]{Alpoge2026}.
Lemma~\ref{lem:an-period-functions} records the form needed here.  We give a
condensed, self-contained proof in the notation of
Section~\ref{sec:analytic-construction}, following the four stages of the
argument in \citep[Sections~3.1--3.4]{Alpoge2026}.

\begin{proof}[Proof of Lemma~\ref{lem:an-period-functions}]
\smallskip
\noindent\emph{The modular lift and its cusp normalization.}
We use the standard facts about the modular group and the modular invariant
collected in
\citep[Chapter~VII, Sections~1.1--1.2 and~3.1--3.3]{Serre1973}.  Let $j$ have its usual
normalization $j(i)=1728$, and set $\mathcal J=1728t\circ\pi$ and
$Z^\circ=\mathfrak H_z\setminus\mathcal J^{-1}\{0,1728\}$.
After the elliptic orbits are removed, $j$ is an unramified modular
covering of $\CC\setminus\{0,1728\}$.  The fundamental group of
$Z^\circ$ is normally generated by small circles around the deleted
points.  Their $\mathcal J$-images are cubes of meridians about zero and fourth
powers of meridians about $1728$.  Since the corresponding modular deck
transformations have orders three and two, respectively, the
covering-space lifting criterion gives a lift
$\tau:Z^\circ\to\mathfrak H$ of $\mathcal J$.  In a bounded-disc model of
$\mathfrak H$, removability and the maximum principle extend the lift
across all deleted points.

For each $g\in\Gamma_{\mathrm{orb}}$, the maps $\tau\circ g$ and $\tau$ lift the
same function.  Thus $\tau(gz)=\phi(g)\tau(z)$ for a homomorphism
$\phi:\Gamma_{\mathrm{orb}}\to\operatorname{PSL}_2(\ZZ)$.  Normalize
$\tau(z_1)=\rho=e^{\pi i/3}$.  The map $\mathcal J$ has order three at
$z_1$, while $j$ has order three at $\rho$, so $\tau-\rho$ has
order one.  Comparing its rotation with
$s_1(g_1z)=e^{-2\pi i/3}s_1(z)$ forces
$\phi(g_1)=\sigma_1$, where
$\sigma_1(\tau)=(\tau-1)/\tau$.
At $z_2$, the function $\mathcal J-1728$ has order four whereas
$j-1728$ has order two at $i$.  Hence $\tau-\tau(z_2)$ has order
two.  In particular $\phi(g_2)\ne1$, for otherwise an invariant germ
would have order divisible by four.  Since the modular group has no
element of order four, $\phi(g_2)$ has order two.

To determine $P=\phi(g_0)$, take $R$ large.  Every component
of $\{|j|>R\}$ is contained in a horoball based at a rational cusp, and
distinct such horoballs are disjoint.  Indeed, this follows from the
standard modular fundamental domain together with
$\Im((a\tau+b)/(c\tau+d))=\Im\tau/|c\tau+d|^2$.
The chosen component $\widetilde U_0$ maps into one such horoball, so $P$
fixes the corresponding rational cusp.  It is not the identity.  Otherwise
$q=e^{2\pi i\tau}$ would descend to the punctured $t_c$-disc and
extend boundedly across zero.  The facts that $j\to\infty$ and
$|q(0)|\ne1$ would give $q(0)=0$, but the single-valued logarithm
$2\pi i\tau$ would force the winding number of $q$ to vanish.  Thus
$P$ is nontrivial parabolic.

Choose its trace-two lift in the form
\[
 P=I+k
 \begin{pmatrix}
 -pr&p^2\\-r^2&pr
 \end{pmatrix},
 \qquad \gcd(p,r)=1,\quad k\ne0.
\]
The element $\phi(g_2)=\sigma_1^{-1}P^{-1}$ has trace zero.  Substitution
gives $-k(p^2-pr+r^2)=1$.
Consequently $k=-1$ and
$(p,r)=\pm(1,0),\pm(0,1),\pm(1,1)$.  Conjugating by a power of $\sigma_1$,
which preserves the normalization at $z_1$, gives
$P(\tau)=\tau-1$ and $\phi(g_2)(\tau)=-1/\tau$.
Hence \eqref{eq:an-tau-laws} holds with the stated cusp normalization.

The Fourier expansion of $j$, together with
$j(\tau)=1728/t_c$, gives $q=e^{2\pi i\tau}=t_cu_1(t_c)$
for a holomorphic unit $u_1$.  After shrinking the cusp disc, choose a
holomorphic logarithm and set
$h(t_c)=(2\pi i)^{-1}\log u_1(t_c)$ and $s=\tau-h(t_c)$.
Then \eqref{eq:an-log-coordinate} holds.  After shrinking $U_0^\times$, the
maps from $\widetilde U_0$ and from a half-plane in the $s$-plane to
$U_0^\times$ are universal covers with the same deck generator; hence
$s$ is a logarithmic coordinate.  In particular,
$\Im\tau\to\infty$ at the cusp.

\smallskip
\noindent\emph{Construction of $\mu$.}
For $\mu$, differences of solutions
of \eqref{eq:an-mu-laws} obey $\nu(g_1z)=-\nu/\tau$ and
$\nu(g_2z)=\nu/\tau$.
The products of these factors around $g_1^3$ and $g_2^4$ are one, so
they extend to a factor of automorphy for $\Gamma_{\mathrm{orb}}$.  Let
$\mathcal L$ be the sheaf on $B$ whose local sections are holomorphic
functions satisfying the resulting equivariance law.  At the cusp, its
sections are required to be holomorphic in $t_c$.  We claim
$\mathcal L\simeq\cO_{\PP^1}(-p_0)$.
To see that $\mathcal L$ is locally free at $p_1,p_2$ and identify it,
one must account for the stabilizers.  On the simply connected cover the
meromorphic function
$F_{\mathcal L}=E_4(\tau)^2E_6(\tau)^{1/2}/
\Delta_{\mathrm{mod}}(\tau)$
is defined because the pullback $E_6(\tau(z))$ has zeros of even order
two over $p_2$.  The modular transformation laws give precisely the
preceding homogeneous factors, up to a possible sign character.  More
precisely, $F_{\mathcal L}(g_1z)=-F_{\mathcal L}(z)/\tau(z)$ and
$F_{\mathcal L}(g_2z)=F_{\mathcal L}(z)/\tau(z)$.
The relation $g_1^3=1$ fixes the sign in the first formula, while comparison
of the leading term in $s_2(g_2z)=-is_2(z)$ fixes the sign in the second.

Upstairs, $F_{\mathcal L}$ has orders two and one at $z_1,z_2$, and
at the cusp it is $t_c^{-1}$ times a unit.  If $\nu$ is a local
holomorphic section, the quotient $\nu/F_{\mathcal L}$ is invariant.  At
$z_1$ its order is at least $-2$ and divisible by three, while at
$z_2$ its order is at least $-1$ and divisible by four.  Thus
$\nu/F_{\mathcal L}$ descends to a function holomorphic at $p_1$ and
$p_2$.  At $p_0$, the
regularity condition is exactly
$\mathcal L_{p_0}=\mathfrak m_{p_0}F_{\mathcal L}$.  Thus
$\mathcal L\simeq\cO_{\PP^1}(-p_0)$.

The affine laws \eqref{eq:an-mu-laws} close after three and four
iterations:
\[
\begin{aligned}
(\tau,\mu)&\xmapsto{g_1}
 \left(\frac{\tau-1}{\tau},\frac{1-\mu}{\tau}\right)
 \xmapsto{g_1}
 \left(\frac{-1}{\tau-1},\frac{\tau-1+\mu}{\tau-1}\right)
 \xmapsto{g_1}(\tau,\mu),\\
(\tau,\mu)&\xmapsto{g_2}
 \left(\frac{-1}{\tau},1+\frac\mu\tau\right)
 \xmapsto{g_2}(\tau,1-\tau-\mu)
 \xmapsto{g_2}
 \left(\frac{-1}{\tau},\frac{1-\mu}{\tau}\right)
 \xmapsto{g_2}(\tau,\mu).
\end{aligned}
\]
Explicit local solutions are $\mu=(2-\tau)/3$ at $p_1$,
$\mu=(1-\tau)/2$ at $p_2$, and $\mu=0$ on $\widetilde U_0$.
At an ordinary point, choose a solution on one sheet and extend it
equivariantly.  The solution sheaf is a torsor under $\mathcal L$, and
$H^0(\PP^1,\cO(-1))=H^1(\PP^1,\cO(-1))=0$.
Thus the local solutions glue uniquely to $\mu$.

\smallskip
\noindent\emph{Construction of $\beta$.}
For $\beta$, put $\varphi_1=2-6(1-\mu)^2/\tau$ and
$\varphi_2=-3-6\mu^2/\tau$.
Substitution in \eqref{eq:an-tau-laws}--\eqref{eq:an-mu-laws} gives
\begin{equation}\label{eq:an-beta-zero-sums}
 \sum_{k=0}^{2}\varphi_1(g_1^kz)=0,
 \qquad
 \sum_{k=0}^{3}\varphi_2(g_2^kz)=0.
\end{equation}
Hence the affine laws \eqref{eq:an-beta-laws} close.  A local solution at
$p_j$ is $m_j^{-1}\sum_{k=0}^{m_j-1}k\varphi_j(g_j^kz)$,
because \eqref{eq:an-beta-zero-sums} makes its increment telescope to
$\varphi_j$.  At the cusp take $\beta=-\tau$, so that
$b_c=\beta+\tau$ is regular.  Differences of solutions form
$\cO_{\PP^1}$, and $H^1(\PP^1,\cO)=0$, so a global $\beta$ exists,
unique up to a constant.

\smallskip
\noindent\emph{Nondegeneracy.}
Finally, direct substitution shows that the function $\mathcal D$ in
\eqref{eq:an-D-negative} is
$\Gamma_{\mathrm{orb}}$-invariant.  It therefore descends to a continuous
function on $B\setminus\{p_0\}$, and at the cusp
$\mathcal D\leq \Im b_c-\Im\tau\to-\infty$.
It is bounded above off a small cusp disc.  Adding a sufficiently negative
imaginary constant to $\beta$ therefore gives
\eqref{eq:an-D-negative} everywhere.
\end{proof}
